\section{Conclusão}

Esta revisão sistematizou 41 estudos sobre frameworks de desenvolvimento de software com IA que deslocam a unidade de trabalho do prompt isolado para um processo estruturado, com artefatos persistentes, papéis agênticos, execução integrada e validação. A tese de que esse deslocamento já ocorreu apoia-se na evidência: a contribuição é descrita como framework em 37 dos 41 estudos, e o diferenciador declarado é a estrutura de processo, não o modelo de linguagem subjacente. Sobre esse corpus, propusemos e testamos uma taxonomia de seis dimensões. Cinco dimensões (especificação, contexto, papéis, execução e validação) recebem suporte empírico forte; portabilidade é fraca e predominantemente diagnóstica, e dois eixos relevantes, o humano-organizacional e o de segurança, emergem fora da taxonomia e recomendam estendê-la.

O campo é jovem, acelera desde 2024 e ainda não tem foro de publicação consolidado, o que se reflete na dispersão de venues e na presença expressiva de preprints. A evidência é desigual, e três eixos de tensão (autonomia versus controle humano, governança pesada versus agilidade leve e especificar antes versus recuperar a especificação do legado) estruturam o debate sem resolvê-lo. Os gargalos práticos mais citados não estão na geração de código, mas em manter o agente ancorado no projeto real e auditável: contexto e rastreabilidade são os requisitos transversais dominantes.

A consequência para a pesquisa futura é direta. O próximo passo do campo é empírico e orientado a processo: benchmarks que avaliem o pipeline e não só a solução, comparações diretas entre frameworks, instrumentos de portabilidade, avaliação longitudinal de governança, uma camada sociotécnica na classificação e avaliação sistemática de segurança. Esta revisão oferece a taxonomia validada e a agenda que organizam esse esforço.
